% Source: https://artofproblemsolving.com/community/c3h1268289_volume_prism
% Stable identity: aops-topic-1268289-q01
Let the edge lengths be $a$, $b$, $c$, with face areas $A=ab$, $B=bc$, $C=ca$.

\begin{asy}
usepackage("amsmath");
size(4.5cm);
defaultpen(fontsize(10pt));
draw((0,0)--(4,0)--(4,3)--(0,3)--cycle);
draw((4,0)--(5.5,1)--(5.5,4)--(4,3)); draw((0,3)--(1.5,4)--(5.5,4));
draw((0,0)--(1.5,1)--(5.5,1),dashed); draw((1.5,1)--(1.5,4),dashed);
label("$a$",(2,0),S); label("$b$",(4.75,0.5),SE); label("$c$",(0,1.5),W);
label("\mbox{}",(min(currentpicture,user=true)+max(currentpicture,user=true))/2,black); // keeps the text after the figure black when LaTeX typesets the labels inline
\end{asy}

Then
\[ABC=(ab)(bc)(ca)=(abc)^2,\]
so the volume is $V=abc=\sqrt{ABC}$.

For areas $6$, $8$, $12$: $V=\sqrt{6\cdot8\cdot12}=\sqrt{576}=\boxed{24}$ (the box is $2\times 3\times 4$).
